% Placeholder skeleton. Replace with a UJ-approved template once chosen.
% Candidate templates:
%   * https://www.overleaf.com/latex/templates/completely-unofficial-jagiellonian-university-thesis-template/fxnssjzrtpfk
%     (Unofficial UJ; supports English and Polish.)
%   * https://www.overleaf.com/latex/templates/wzor-isi-uj/ddszfvqhkwmt
%     (UJ ISI; less appropriate — different faculty.)
%   * Ask supervisor Lech Duraj if WMiI UJ has an official template.
\documentclass[12pt,a4paper]{report}

\usepackage[T1]{fontenc}
\usepackage[utf8]{inputenc}
\usepackage[english]{babel}
\usepackage{lmodern}
\usepackage{microtype}
\usepackage{amsmath, amssymb, amsthm}
\usepackage{algorithm, algpseudocode}
\usepackage{graphicx}
\usepackage{booktabs}
\usepackage{hyperref}
\usepackage{cleveref}
\usepackage[a4paper, margin=2.5cm]{geometry}

\newtheorem{theorem}{Theorem}[chapter]
\newtheorem{lemma}[theorem]{Lemma}
\newtheorem{definition}[theorem]{Definition}

\title{Implementation and Performance Analysis of a Polynomial-Time Algorithm
       for the Next-to-Shortest Path Problem in Directed Weighted Graphs}
\author{Ihar Maroz \\ Supervisor: Lech Duraj \\ Jagiellonian University}
\date{2026}

\begin{document}

\maketitle

\begin{abstract}
TODO: write abstract once experiments are in. Briefly: we implement the
Chen--Wein--Zhang (2025) polynomial-time algorithm for the next-to-shortest
path problem on positively weighted directed graphs, the first polynomial
algorithm for this problem after roughly thirty years of being open, and
empirically evaluate its practical performance against a brute-force oracle
and Yen's $k$-shortest-simple-paths heuristic with $k=2$.
\end{abstract}

\tableofcontents

\chapter{Introduction}

% TODO: motivate the next-to-shortest path problem; state the historical openness
% of the directed positively-weighted case; preview the contribution of this work
% (first empirical study of CWZ).

\chapter{Literature Review}
\label{ch:litreview}

\section{Shortest paths}

Single-source shortest paths in directed graphs with non-negative weights are
classically solved by Dijkstra's algorithm~\cite{Dijkstra1959}, which we use
throughout in its binary-heap form, running in $O((|V|+|E|)\log|V|)$ time.

\section{$k$-shortest simple paths}

Yen's algorithm~\cite{Yen1971} enumerates the $k$ shortest simple $s\to t$
paths in non-decreasing order of weight using $k$ Dijkstra-style searches per
``spur'' along the previously-found paths. Its runtime is
$O(kn(m+n\log n))$, which is polynomial in $k$ but Yen's algorithm is used
here in a slightly non-standard way: we invoke it as an NSP solver by
enumerating paths until one is found whose weight strictly exceeds
$d_G(s,t)$. If the graph has $r$ distinct simple shortest paths from $s$ to
$t$, this requires examining at least $r$ paths before finding one with
strictly greater cost, and $r$ can be exponential in $|V|$. In favorable
inputs Yen is extremely fast, but on adversarial graphs with many alternate
shortest paths it degenerates.

Eppstein~\cite{Eppstein1998} gives an $O((m + n\log n) + k)$ algorithm for
$k$-shortest paths that need not be simple. NSP requires simplicity and so
Eppstein's algorithm does not directly apply.

\section{Next-to-shortest path: positive results}

The NSP problem was introduced in 1996, and Lalgudi and
Papaefthymiou~\cite{LalgudiPapaefthymiou1997} soon after showed it is
NP-complete in directed graphs with non-negative weights (in
particular, when zero weights are permitted). The undirected case admits a
polynomial-time algorithm: Krasikov and Noble~\cite{KrasikovNoble2004}
gave the first, running in $O(n^3 m)$ time, and a sequence of improvements
brought this down to $O(n^3)$, then $O(n^2)$, and finally to
$O(n\log n + m)$ by Wu~\cite{Wu2010}. For \emph{planar} directed graphs
with positive weights, Wu and Wang~\cite{WuWang2015} gave a
polynomial-time algorithm; they also established structural properties of
solutions on general positively weighted digraphs, but no further progress
followed until Chen, Wein and Zhang.

For \emph{general} directed graphs with strictly positive weights, the
problem remained open for nearly three decades. Chen, Wein, and Zhang
\cite{ChenWeinZhang2025} resolved it in November~2025 with an
$O\bigl(|V|^{4}|E|^{3}\log|V|\bigr)$-time algorithm. The algorithm proceeds
by two reductions -- general digraph to $(s,t)$-straight to $(s,t)$-layered
-- followed by a structural enumeration over $6$-tuples of vertices and
edges that, together with a $2$-vertex-disjoint-paths subroutine in a DAG,
suffice to find any NSP if one exists.

\section{Related primitives}

The CWZ algorithm depends critically on a subroutine for finding two
vertex-disjoint paths between prescribed source-sink pairs in a directed
acyclic graph. Fortune, Hopcroft, and Wyllie~\cite{FortuneHopcroftWyllie1980}
gave the first polynomial-time algorithm for the $k$-vertex-disjoint-paths
problem in DAGs for fixed $k$. Tholey~\cite{Tholey2012} subsequently gave a
linear-time algorithm for the special case $k=2$. The CWZ paper cites Tholey
for the $O(|V|+|E|)$ bound on the inner subroutine. In the present work we
implement the simpler Fortune--Hopcroft--Wyllie algorithm for two paths,
which runs in $O(|V|^{2}+|V||E|)$ time (Section~\ref{sec:twovdp}), so our
implementation does not attain the paper's bound.

\chapter{The Chen--Wein--Zhang Algorithm}

% TODO: three reductions; (s,t)-straight and (s,t)-layered graphs; back-edge
% decomposition; pair of disjoint forward paths; Lemma 5.3 (statement and
% proof sketch); 2-VDP-in-DAG subroutine (Tholey 2012); overall complexity
% $O(|V|^4 |E|^3 \log |V|)$.

\chapter{Implementation}
\label{ch:implementation}

\section{Language and build system}

The implementation is in C++20, compiled with GCC~14.2 on Linux (WSL2).
CMake~$\ge 3.20$ orchestrates the build. We deliberately keep external
dependencies minimal: GoogleTest is fetched at configure time via
\texttt{FetchContent}, the Python plotting tooling needs only
\texttt{pandas}, \texttt{matplotlib}, and \texttt{numpy}, and nothing
else is required on the path of the C++ binaries themselves. The choice
of C++ over Python was driven by the high asymptotic complexity of the
CWZ algorithm: even an $O(|V|^{4}|E|^{3})$ inner loop must run on graphs
with $|V|\approx 30$--$50$ to give meaningful empirical data, and
Python's interpretation overhead would limit the reachable size
considerably.

Edge weights are stored as 64-bit signed integers (\texttt{Weight =
std::int64\_t}). The algorithm's edge classification compares
\textit{exact} equalities like $d_S[u]+w(u,v)=d_S[v]$, which
floating-point arithmetic would render unreliable. Generators emit
weights uniformly drawn from $[1, 10]$.

\section{Module layout}

The source tree is split into small modules under \texttt{src/}:
\texttt{graph/} (a directed graph with adjacency lists indexed by stable
edge IDs and both out- and in-neighbour lists), \texttt{shortest\_path/}
(Dijkstra forward and reverse), \texttt{two\_vdp/} (2-VDP-in-DAG and a
brute-force oracle for it), \texttt{cwz/} (reductions and the layered
NSP algorithm), \texttt{baseline/} (brute force and Yen), and
\texttt{cli/} (a thin \texttt{nsp-cli} wrapper). Graph generators live
under \texttt{generators/}; unit tests under \texttt{tests/unit/};
correctness fuzzers under \texttt{tests/correctness/}; benchmarks and
plotting under \texttt{bench/}.

\section{Key data structures}

The \texttt{Graph} class assigns a stable identifier to each edge at
insertion time and exposes both \texttt{out\_edges(v)} and
\texttt{in\_edges(v)} lists of edge IDs. Maintaining the reverse
adjacency in parallel makes reverse Dijkstra a single pass over in-edges
rather than a transpose-then-search operation, and it lets the layered
algorithm classify each edge as forward/back-edge once and store the
result on the side.

For the residual subgraph $G[V\setminus U]$ used inside the layered
algorithm we never materialize the subgraph; instead we run a
vertex-masked Dijkstra that skips any edge whose endpoint lies in
$U$. This avoids allocating an $O(|V|+|E|)$ subgraph for each of the
many $(A,B)$ candidates explored.

\section{Paper-faithful pipeline}
\label{sec:pipeline}

The top-level \texttt{cwz\_nsp(g, s, t)} runs the paper's three-stage
pipeline:
\begin{enumerate}
  \item \texttt{reduce\_to\_straight(g, s, t)} (the paper's
    \textsc{NextSP}) folds every off-shortest-path vertex, replacing it
    with shortcut edges between its in- and out-neighbours, leaving an
    $(s,t)$-straight graph;
  \item \texttt{reduce\_to\_layered} (the paper's
    \textsc{NextSP-Straight}) turns the straight graph into a
    \emph{strictly} $(s,t)$-layered one and, as a by-product, records a
    set of candidate NSP paths (see below);
  \item \texttt{cwz\_nsp\_layered} (the paper's \textsc{NextSP-Layered})
    runs the Lemma~5.3 $6$-tuple enumeration on the strictly layered
    graph;
  \item the answer is the cheapest among the recorded candidates and the
    layered result, mapped back to original-graph edges via a
    \texttt{path\_expand[]} table that each reduction maintains.
\end{enumerate}

\paragraph{\textsc{NextSP-Straight} removes violating back-edges and
records candidates.} The subtle and important part is step~2. The
paper's $(s,t)$-layered definition (Definition~4.1) requires that every
back-edge go \emph{strictly backward} in layer. An $(s,t)$-straight
graph can contain back-edges that go forward (or sideways) in distance
-- an edge $(u,v)$ with $d_G(s,u)\le d_G(s,v)$ but
$d_G(s,u)+w(u,v)>d_G(s,v)$. \textsc{NextSP-Straight} \emph{removes} each
such edge, but first records the candidate path
$P_{s\to u}\circ(u,v)\circ P_{v\to t}$ (shortest forward paths to $u$
and from $v$), which the paper proves is always a valid simple
not-shortest path. Forward edges that skip layers are subdivided through
intermediate-distance vertices; strictly-backward back-edges are kept
unchanged. The result is a graph in which \emph{every} back-edge is
strictly backward -- exactly the structure Lemma~5.1 guarantees for an
$(s,t)$-layered graph, and exactly the structure the layered enumeration
relies on. Our implementation (\texttt{reduce\_to\_layered}) records the
removed-edge candidates in \texttt{ReductionResult::candidates}, and
\texttt{cwz\_nsp} folds them into the final minimum.

\section{The layered NSP algorithm}
\label{sec:layered}

\texttt{cwz\_nsp\_layered} is a direct transcription of
\textsc{NextSP-Layered} (paper Section~5). It enumerates every
$6$-tuple $(A,B,X',X,Y',Y)$ with $A,B$ incident to back-edges,
$\lambda(X')=\lambda(Y')=\lambda(X)-1=\lambda(Y)-1$, and (implied by the
layer constraints) $d(A)>d(B)$. For each tuple it finds the pair of
vertex-disjoint forward paths $P_1=s\to X'\to X\to A$ and
$P_2=B\to Y'\to Y\to t$ via two calls to the 2-VDP-in-DAG subroutine --
one for the layers below the barrier ($s\to X'$, $B\to Y'$) and one for
the layers above ($X\to A$, $Y\to t$), which are automatically disjoint
because forward paths are monotone in layer. It then completes the
candidate with a shortest $A\to B$ path in the residual graph
$G[(V\setminus(V(P_1)\cup V(P_2)))\cup\{A,B\}]$ and keeps the cheapest
result. The Key Lemma (Lemma~5.3) guarantees this finds an NSP whenever
one exists.

The boundary cases $A=X$ (the upper part of $P_1$ is empty) and $B=Y'$
(the lower part of $P_2$ is empty) are common on small graphs and must
be admitted: the layer constraint on $B$ is $\lambda(B)\le L$, not
$\lambda(B)<L$, and infeasible tuples are rejected by the 2-VDP calls
rather than excluded up front. Tightening either of these silently loses
NSPs whose first back-edge sits exactly one layer above its last.

\section{The 2-VDP subroutine}
\label{sec:twovdp}

Given a DAG $D=(V,E)$ and vertices $s_1,t_1,s_2,t_2$, the subroutine
must find paths $P_1\colon s_1\to t_1$ and $P_2\colon s_2\to t_2$ with
$V(P_1)\cap V(P_2)=\emptyset$, or report that no such pair exists. The
layered algorithm calls it on the subgraph of forward edges. That
subgraph is acyclic: a forward edge $(u,v)$ satisfies
$d(s,v)=d(s,u)+w(u,v)>d(s,u)$, since all weights are positive.

We implement the algorithm of Fortune, Hopcroft and
Wyllie~\cite{FortuneHopcroftWyllie1980} for two paths
(\texttt{two\_vdp\_in\_dag}). It runs in $O(|V|^{2}+|V||E|)$ time, which
is slower than Tholey's $O(|V|+|E|)$ algorithm~\cite{Tholey2012} but
much shorter to implement.

\paragraph{Algorithm.} If $s_1=s_2$, $t_1=t_2$, $s_1=t_2$ or $s_2=t_1$,
some vertex would have to lie on both paths, and the answer is
negative. Otherwise fix a topological order of $D$ and let $\pi(x)$ be
the position of vertex $x$. A \emph{state} is a pair $(u,v)$ of distinct
vertices, meaning that $P_1$ has been built from $s_1$ up to $u$ and
$P_2$ from $s_2$ up to $v$. In a state $(u,v)\neq(t_1,t_2)$ exactly one
of the two paths is extended:
\begin{itemize}
  \item $P_1$ is extended if $u\neq t_1$ and either $v=t_2$ or
    $\pi(u)<\pi(v)$; this gives the moves $(u,v)\to(w,v)$ for every edge
    $(u,w)\in E$ with $w\neq v$;
  \item otherwise $P_2$ is extended, giving the moves $(u,v)\to(u,w)$
    for every edge $(v,w)\in E$ with $w\neq u$.
\end{itemize}
In words, the path whose last vertex comes earlier in the topological
order is extended, and a path that has reached its target is never
extended again. If $P_2$ is extended then $v\neq t_2$: either $u\neq
t_1$ and the rule forces it, or $u=t_1$ and then $v\neq t_2$ because the
state is not $(t_1,t_2)$.

The moves form a directed graph $H$ on at most $|V|(|V|-1)$ states. The
algorithm runs a depth-first search in $H$ from $(s_1,s_2)$ and stores,
for every state, the edge of $D$ through which the state was first
reached. If $(t_1,t_2)$ is reached, the paths are recovered by following
the stored edges back to $(s_1,s_2)$. Each stored edge ends in exactly
one of the two current vertices, since $u\neq v$, and that tells us
whether it belongs to $P_1$ or to $P_2$. If the search ends without
reaching $(t_1,t_2)$, the answer is negative.

\paragraph{Correctness.} A walk in $H$ starting at $(s_1,s_2)$ records
two sequences of vertices of $D$, one for each path, by listing the
vertices added by the moves.

\begin{lemma}\label{lem:vdp-sound}
For every walk in $H$ from $(s_1,s_2)$ to a state $(u,v)$, the recorded
sequences are vertex-disjoint paths $P_1\colon s_1\to u$ and
$P_2\colon s_2\to v$.
\end{lemma}

\begin{proof}
Induction on the length of the walk. Before any move is made, the state
is $(s_1,s_2)$, $P_1$ consists of the single vertex $s_1$ and $P_2$ of
the single vertex $s_2$. They are disjoint because $s_1\neq s_2$. Suppose the claim holds in $(u,v)$ and
the next move extends $P_1$ along an edge $(u,w)$. Since $D$ is acyclic,
$\pi(w)>\pi(u)\ge\pi(x)$ for every $x\in V(P_1)$, so $P_1$ remains a
path. It remains to show $w\notin V(P_2)$. The move requires $w\neq v$.
Let $x$ be a vertex of $P_2$ other than $v$. Then $P_2$ has been
extended at least once; consider its last extension, a move
$(u',v')\to(u',v)$. The rule chose $P_2$ in state $(u',v')$, so either
$u'=t_1$, or $v'\neq t_2$ and $\pi(v')<\pi(u')$. If $u'=t_1$, then $P_1$
was never extended afterwards and $u=t_1$, which contradicts the fact
that $P_1$ is extended now. Hence $\pi(v')<\pi(u')$. The vertex $x$ lies
on $P_2$ no later than $v'$ and $u'$ lies on $P_1$ no later than $u$, so
\[
  \pi(x)\le\pi(v')<\pi(u')\le\pi(u)<\pi(w),
\]
and $x\neq w$. The case in which $P_2$ is extended is symmetric: the
last extension of $P_1$ took place in a state $(u',v')$ with $u'\neq
t_1$ and either $v'=t_2$ or $\pi(u')<\pi(v')$; the first option would
mean $v=t_2$, but $P_2$ is being extended, so
$\pi(x)\le\pi(u')<\pi(v')\le\pi(v)<\pi(w)$ for every $x\in V(P_1)$
other than $u$.
\end{proof}

\begin{lemma}\label{lem:vdp-complete}
If there are vertex-disjoint paths $Q_1\colon s_1\to t_1$ and
$Q_2\colon s_2\to t_2$ in $D$, then $(t_1,t_2)$ is reachable from
$(s_1,s_2)$ in $H$.
\end{lemma}

\begin{proof}
We build a walk in $H$ that keeps $u$ on $Q_1$ and $v$ on $Q_2$. This
holds in $(s_1,s_2)$. In a state $(u,v)\neq(t_1,t_2)$ the rule selects a
path that has not reached its target, as observed above. Move its last
vertex to the next vertex $w$ of $Q_1$ or $Q_2$ respectively. Since
$Q_1$ and $Q_2$ are disjoint, $w$ differs from the last vertex of the
other path, so this is a move of $H$, and the invariant still holds.
Every move advances along $Q_1$ or $Q_2$, so after
$|E(Q_1)|+|E(Q_2)|$ moves the walk is in $(t_1,t_2)$.
\end{proof}

The depth-first search visits every state reachable from $(s_1,s_2)$,
so by Lemma~\ref{lem:vdp-complete} it finds $(t_1,t_2)$ whenever a
solution exists. The stored edges describe one walk from $(s_1,s_2)$ to
$(t_1,t_2)$, and by Lemma~\ref{lem:vdp-sound} the paths recovered from
it are vertex-disjoint. Hence \texttt{two\_vdp\_in\_dag} returns a pair
of paths if and only if one exists, and every returned pair is correct.

\paragraph{Running time.} The topological order is computed with
Kahn's algorithm in $O(|V|+|E|)$ time; if the input contains a cycle,
the order is incomplete and the function throws instead of answering.
Each state is expanded at most once. Expanding $(u,v)$ scans the
out-edges of one of $u$ and $v$, so the total work of the search is at
most
\[
  \sum_{u\in V}\sum_{v\in V}\bigl(\deg^{+}(u)+\deg^{+}(v)\bigr)
  = 2|V||E|.
\]
Together with the $|V|^{2}$-entry array of stored edges this gives
$O(|V|^{2}+|V||E|)$ time and $O(|V|^{2})$ memory per call.

\section{Testing strategy}
\label{sec:correctness}

The project ships:
\begin{itemize}
  \item $110$ unit tests covering the graph data structure, both
    Dijkstras, 2-VDP ($54$ tests), both
    reductions (including that \textsc{NextSP-Straight} removes
    forward/sideways back-edges and records the right candidates), the
    layered NSP algorithm on hand-built examples, both baselines, and
    the end-to-end \texttt{cwz\_nsp} pipeline. All pass.
  \item A brute-force 2-VDP solver
    (\texttt{src/two\_vdp/brute\_force\_two\_vdp}) used as an oracle for
    \texttt{two\_vdp\_in\_dag}. It enumerates all simple $s_1\to t_1$
    and $s_2\to t_2$ paths by depth-first search, stores the vertex set
    of each path as a $64$-bit mask, and reports a solution if and only
    if some $s_1\to t_1$ mask and some $s_2\to t_2$ mask have an empty
    intersection. It takes exponential time and is limited to $64$
    vertices, but it shares no code with the Fortune--Hopcroft--Wyllie
    implementation and does not rely on acyclicity. The unit tests
    compare the two on every terminal quadruple $(s_1,t_1,s_2,t_2)$ of
    all DAGs on $4$ labelled vertices ($393{,}216$ queries), of all
    $1{,}024$ edge subsets of a $5$-vertex transitive tournament under a
    random relabelling ($640{,}000$ queries), and of $400$ random DAGs
    with up to $7$ vertices and occasional parallel edges ($266{,}533$
    queries). Every pair of paths returned by \texttt{two\_vdp\_in\_dag}
    is also checked to consist of two vertex-disjoint paths with the
    right endpoints. The two solvers agree on all $1{,}299{,}749$
    queries. Hand-built tests cover the degenerate cases $s_i=t_i$,
    parallel edges, rejection of cyclic inputs, and instances in which
    only the crossed pairing $s_1\to t_2$, $s_2\to t_1$ is routable.
  \item A property-based fuzz \texttt{fuzz\_brute\_vs\_yen} that
    confirms brute force and Yen's NSP heuristic agree on cost across
    thousands of random small graphs (mismatches would indicate a bug
    in either side).
  \item A differential fuzz \texttt{fuzz\_cwz\_vs\_brute} that confirms
    the full CWZ pipeline matches brute force. We measured \emph{zero}
    mismatches across $34{,}000$ random trials drawn from $35$ distinct
    RNG seeds: $|V|\le 10$ over $20$ seeds ($20{,}000$ trials),
    $|V|\le 12$ over $10$ seeds ($10{,}000$ trials), and $|V|\le 14$
    over $5$ seeds ($4{,}000$ trials). Of these, $23{,}624$ instances
    had an NSP and the remaining $10{,}376$ had none; both algorithms
    agreed on which case held every time. The seeds are consecutive
    integers ($1001$--$1020$, $2001$--$2010$, $3001$--$3005$), so the
    campaign is reproducible verbatim.
\end{itemize}

The combined test suite establishes correctness on small graphs;
agreement with brute force on all $1200$ benchmark instances
(Section~\ref{sec:agreement}) extends the evidence to larger graphs.

\section{The CLI}

\texttt{nsp-cli} reads a graph from a DIMACS-style text file (first line
$n\;m$, then $m$ lines of $u\;v\;w$) and runs the selected algorithm.
Source and sink default to vertex $0$ and vertex $n-1$. The output is
one \texttt{nsp\_cost} line and a list of edges describing the recovered
NSP.

\chapter{Experiments}
\label{ch:experiments}

\section{Methodology}
\label{sec:methodology}

We benchmark three algorithms -- brute-force NSP, Yen-NSP, and our
implementation of the Chen--Wein--Zhang (CWZ) algorithm -- on graphs
drawn from five generator families:
\begin{itemize}
  \item Erd\H{o}s--R\'enyi random digraphs $G_{n,p}$ with $p=0.30$;
  \item random DAGs (vertices in topological order; each forward edge
    present with probability $p=0.30$);
  \item layered digraphs (two endpoint vertices $s$, $t$ plus $L$ interior
    layers of width $W$, with forward-edge probability $p_f=0.40$ and
    back-edge probability $p_b=0.10$). The name refers to this generator's
    layer topology and \emph{not} to the $(s,t)$-layered property of
    Definition~4.1: these instances do not satisfy it -- a typical one has
    most of its vertices off every shortest $s\to t$ path and most of its
    back-edges not strictly backward -- so both reductions run in full on
    them, exactly as on the other families;
  \item grid digraphs (rows-by-columns lattice with rightward and downward
    edges only);
  \item \emph{diamond-chain} digraphs, a deliberately adversarial family
    described in Section~\ref{sec:adversarial}.
\end{itemize}

For each family we sweep $|V|\in\{6,8,10,12,16,20,25,30,40,50,60,70\}$, drawing
$20$ instances per size with deterministic seeds. Any instance in which $t$
is unreachable from $s$ is redrawn, so every measurement is taken on a graph
that actually poses the problem (Section~\ref{sec:existence}). All three
algorithms run across the whole sweep. The brute-force oracle is a branch-and-bound search
rather than an exhaustive enumeration: a reverse Dijkstra supplies the
admissible lower bound $d(u,t)$, and any prefix with
$\text{cost}+d(u,t)\ge\text{best}$ is pruned. That bound is what makes an
exponential-in-the-worst-case oracle usable across the sweep; it collapses,
by construction, only on graphs with very many equal-cost shortest paths,
where every shortest path must still be walked. Each algorithm call is
wrapped in a $30$-second wall-clock
budget; calls that exceed it are recorded as ``budget exceeded'' and
excluded from the timing aggregates. Edge weights are integers drawn
uniformly from $[1,10]$. All three algorithms are implemented in C++20,
compiled at \texttt{-O3}, and timed with wall-clock measurements taken
\emph{around the algorithm call only} (graph generation excluded). The
benchmark binary (\texttt{build/bench}) writes one CSV row per
$\langle\text{algorithm},\text{family},|V|,\text{trial}\rangle$ and the
plots are produced by \texttt{bench/plots/make\_plots.py}.

For each $(\text{family},|V|)$ cell we report the \emph{median} runtime
across the $20$ trials. Plots additionally show the inter-quartile range
($25^{\text{th}}$--$75^{\text{th}}$ percentile) as a shaded band around
the median, which is robust to single-trial outliers and to the
heavy-tailed runtimes that appear on dense Erd\H{o}s--R\'enyi instances.

\section{Runtime by family and size}
\label{sec:runtime-by-family}

Figure~\ref{fig:runtime} shows the runtime breakdown by family, with one
sub-panel per family. The vertical axis is logarithmic and the shaded
band around each line is the IQR. Three observations stand out.

\paragraph{All three algorithms are fast on the benign families.}
Even at $|V|=50$, CWZ on the densest family (Erd\H{o}s--R\'enyi) has a
median of $208$~ms at $|V|=70$ ($59$~ms at $|V|=50$). Brute force, despite its exponential worst case, is
the \emph{fastest} of the three on the four benign families: its lower-bound
pruning restricts the search to paths that could still beat the best
not-shortest path found so far, and on these inputs that is a handful of
short prefixes. Yen is likewise fast throughout, finishing in tens of
microseconds at $|V|=50$.

\paragraph{CWZ's runtime is family-sensitive.} On Erd\H{o}s--R\'enyi
graphs the median grows from $51$~\textmu s at $|V|=10$ to $208$~ms at
$|V|=70$, a four-orders-of-magnitude climb driven by the number of
back-edges that the inner enumeration must consider. On the random-DAG
family the same sweep climbs only from $11$~\textmu s to $3.4$~ms because
DAGs lack the dense back-edge structure of $G_{n,p}$. Layered and grid
graphs sit in between.

\paragraph{Yen dominates on the four random families.} For the input
distributions we tested, Yen-NSP terminates orders of magnitude faster
than CWZ. This is the expected behaviour on graphs where the number of
distinct simple shortest paths is small: Yen only has to step past a
handful of paths before finding one with strictly greater weight, while
CWZ pays its polynomial price uniformly. The interest of CWZ is precisely
in \emph{worst-case} polynomiality -- it is the only algorithm with a
fixed polynomial bound regardless of how many shortest paths the graph
contains. The next subsection shows that worst-case behaviour
empirically.

\paragraph{The diamond-chain panel inverts the picture.} The rightmost
sub-panel of Figure~\ref{fig:runtime} shows the same axes for the
diamond-chain family. Yen's curve climbs steeply -- from microseconds
at $|V|=10$ to $3.7$~s at $|V|=70$ -- while the CWZ curve stays
microsecond-flat at $51$~\textmu s. This is the empirical face of the
$2^k$ blow-up Yen must endure on graphs with exponentially many shortest
paths, against which CWZ's polynomial bound is the only protection.

\begin{figure}
  \centering
  \includegraphics[width=\textwidth]{../results/runtime_by_family.pdf}
  \caption{Median runtime versus $|V|$ across families (log $y$-axis), with
    inter-quartile range shaded. Each family is a sub-panel; the
    diamond-chain panel inverts the brute-vs-Yen-vs-CWZ ordering visible in
    the other four, with both exponential algorithms climbing past CWZ while
    CWZ stays flat. Twenty trials per $(\text{family},|V|)$ cell. Points
    sitting exactly on the $1$~\textmu s line are below our wall-clock timer
    resolution and are clamped there for display, so a flat segment at that
    level means ``too fast to measure'', not constant time.}
  \label{fig:runtime}
\end{figure}

\section{CWZ scaling}

Figure~\ref{fig:scaling} shows the CWZ median runtime alone on a log-log
scale, with all five families overlaid. Fitting a power law to the median
runtime of each family separately gives exponents that vary widely with
graph structure:

\begin{center}
\begin{tabular}{lc}
\toprule
family & log-log slope \\
\midrule
diamond chain  & $1.29$ \\
grid           & $2.28$ \\
layered        & $2.90$ \\
random DAG     & $2.90$ \\
Erd\H{o}s--R\'enyi & $4.34$ \\
\bottomrule
\end{tabular}
\end{center}

The spread tracks edge density: the sparse structured families sit near
$|V|^{1}$--$|V|^{3}$, while dense Erd\H{o}s--R\'enyi graphs
($|E|=\Theta(|V|^{2})$ at $p=0.30$) reach $\approx |V|^{4.3}$. The fits use each instance's
\emph{actual} vertex count, not the sweep target, which differ by up to
$30\%$ for the grid and diamond-chain families. Every
family remains far below the paper's worst-case bound -- for
$|E|=\Theta(|V|)$ the $|V|^{4}|E|^{3}$ bound is a slope of $\approx 7$ in
$\log|V|$, and for the dense Erd\H{o}s--R\'enyi regime it is
$\approx 10$. This indicates that the inner enumeration of our
implementation prunes aggressively in practice on these inputs, and that
the observed exponent is a property of the input family rather than a
single characteristic number.

\begin{figure}
  \centering
  \includegraphics[width=0.85\textwidth]{../results/cwz_scaling.pdf}
  \caption{CWZ median runtime on a log-log scale, all five families.}
  \label{fig:scaling}
\end{figure}

\section{Agreement between algorithms}
\label{sec:agreement}

The strongest check available is against the exact oracle. Because the
brute-force search now runs across the entire sweep, every one of the
$1200$ benchmark instances can be compared against it directly:
\textbf{CWZ agrees with brute force on all $1200$}, spanning all five
families and $|V|$ up to $81$. This extends the correctness evidence well
beyond the $|V|\le 14$ reach of the fuzz harness of
Chapter~\ref{ch:implementation}.

Against Yen the picture is nearly as clean. The script
\texttt{bench/plots/make\_plots.py} emits
\texttt{results/agreement\_summary.csv} which records, per
$(\text{family},|V|)$ cell, how many of the $20$ trials had matching NSP
costs. Agreement is $100\%$ on the four benign
families across the whole sweep, and on the diamond-chain family up to
$|V|=30$.

The only cells with less than full agreement are the two largest
diamond-chain sizes, $|V|=40$ and $|V|=50$ -- and there it is
\emph{Yen}, not CWZ, that is wrong. On those instances the number of
distinct shortest paths ($2^{k}$ with $k\ge 13$) exceeds Yen's
enumeration cap $k_{\max}=5000$, so Yen exhausts its budget and reports
``no NSP'' even though one exists. CWZ returns the correct NSP (cost
$2k+1$) in tens of microseconds on the same instances. This is a sharper
form of the adversarial story of Section~\ref{sec:adversarial}: a
capped Yen does not merely run slowly on graphs with exponentially many
shortest paths, it silently returns the wrong answer.

This agreement check is complementary to the fuzz harness of
Chapter~\ref{ch:implementation}: not only does our CWZ implementation
match the brute-force oracle on small inputs, it matches Yen on every
benign instance and \emph{beats} Yen (matching brute force) on the
adversarial instances where Yen's cap fails.

\begin{table}
  \centering
  \begin{tabular}{lrrr}
  \toprule
  family & $|V|$ range & trials per cell & agree \\
  \midrule
  Erd\H{o}s--R\'enyi & 6--50 & 20 / size & 100\% \\
  random DAG        & 6--50 & 20 / size & 100\% \\
  layered           & 6--50 & 20 / size & 100\% \\
  grid              & 6--50 & 20 / size & 100\% \\
  diamond chain     & 6--30 & 20 / size & 100\% \\
  diamond chain     & 40--50 & 20 / size & Yen capped$^{*}$ \\
  \bottomrule
  \end{tabular}
  \caption{CWZ vs Yen agreement on benchmark instances. The summary CSV
    \texttt{results/agreement\_summary.csv} carries the per-cell
    breakdown. $^{*}$On diamond chains at $|V|\ge 40$ Yen exceeds its
    $k_{\max}=5000$ enumeration cap and reports no NSP; CWZ returns the
    correct NSP, matching brute force, in tens of microseconds.}
  \label{tab:agreement}
\end{table}

\section{NSP existence}
\label{sec:existence}

A practical caveat in interpreting Figure~\ref{fig:runtime}: on some of the
smallest generated instances no NSP exists at all. Figure~\ref{fig:existence}
reports the fraction on which one does, as determined by the brute-force
oracle.

Two effects were at work here, and only one of them is interesting. The first
was an artefact: a graph in which $t$ is simply unreachable from $s$ has no
shortest path, so the question of a next-to-shortest one is vacuous. Sparse
families produced such graphs often at small $|V|$ -- $13/20$
Erd\H{o}s--R\'enyi instances at $|V|=6$ -- and in the layered family the cause
was structural rather than statistical: with $\text{width}=2$ and
$p_f=0.40$, a pair of consecutive layers receives no forward edge at all with
probability $(1-p_f)^{\text{width}^2}=0.6^{4}\approx 13\%$, severing the
graph. The generator now guarantees at least one forward edge between every
consecutive pair of layers, and the benchmark redraws any instance in which
$t$ is unreachable from $s$. Every one of the $1200$ instances reported here
is therefore genuinely solvable.

What remains is real. Of those $1200$ instances, $1122$ have an NSP and $78$ do
not, every one of the latter at $|V|\le 16$. These are graphs with exactly one
simple $s\to t$ path -- which is then the shortest path by default, leaving
nothing to be next-to-shortest -- or, in a handful of cases, several paths
whose costs all tie at the minimum, so that none of them is
\emph{not}-shortest. Both are legitimate NSP instances on which the correct
answer is ``none'', and all three algorithms are required to say so; we keep
them deliberately, as they are the only coverage the benchmark has of that
code path. The effect disappears as soon as the graphs are large enough to
admit several distinct-cost routes: every family is at $100\%$ by $|V|=20$.

It is worth noting how this figure was obtained. An earlier version plotted
the fraction on which \emph{Yen} reported an NSP, which made the curve for
the diamond-chain family collapse to zero at $|V|\ge 40$ -- on the one
family that provably always has an NSP. What it was really measuring was
Yen exhausting its $k_{\max}$ enumeration cap
(Section~\ref{sec:agreement}). Existence is a property of the instance,
not of the solver, so it must be read off an exact oracle.

\begin{figure}
  \centering
  \includegraphics[width=0.75\textwidth]{../results/nsp_existence.pdf}
  \caption{Fraction of generated instances on which an NSP exists, as
    determined by the brute-force oracle.}
  \label{fig:existence}
\end{figure}

\section{Adversarial inputs: where CWZ wins}
\label{sec:adversarial}

The diamond-chain family in Section~\ref{sec:methodology} (which already appears as the
right-most panel of Figure~\ref{fig:runtime}) is repeated here with a
dedicated sweep over $k$ to make the $2^{k}$ scaling visible at higher
resolution.

\paragraph{Diamond-chain construction.} For a parameter $k$ we build a
layered digraph as follows. Vertex $0$ is $s$. For each $i\in\{1,\dots,k\}$
we add three vertices: $\text{top}_i$, $\text{bot}_i$, and $\text{merge}_i$;
$\text{merge}_k$ is $t$. Set $\text{merge}_0\equiv s$. For each $i$, add
four unit-weight edges:
$\text{merge}_{i-1}\to\text{top}_i$, $\text{merge}_{i-1}\to\text{bot}_i$,
$\text{top}_i\to\text{merge}_i$, and $\text{bot}_i\to\text{merge}_i$.
Finally add a direct edge $s\to t$ of weight $2k+1$.

The construction has $1+3k$ vertices, $4k+1$ edges, and exactly $2^{k}$
distinct simple shortest $s\to t$ paths (each of weight $2k$). The unique
NSP is the direct $s\to t$ shortcut of weight $2k+1$.

\paragraph{Yen's exponential blowup.} On this family Yen must enumerate
all $2^{k}$ shortest simple paths before reaching the NSP.
Table~\ref{tab:adv} shows the measured runtimes from the dedicated
\texttt{build/bench\_adversarial} sweep; Figure~\ref{fig:adversarial}
plots them.

\begin{table}
  \centering
  \begin{tabular}{rrrrr}
  \toprule
  $k$ & $|V|$ & $|E|$ & Yen (s) & CWZ (s) \\
  \midrule
  10 & 31 & 41 & 0.023 & 0.000014 \\
  12 & 37 & 49 & 0.332 & 0.000016 \\
  14 & 43 & 57 & 10.04 & 0.000018 \\
  15 & 46 & 61 & 56.18 & 0.000022 \\
  20 & 61 & 81 & --   & 0.000025 \\
  24 & 73 & 97 & --   & 0.000031 \\
  \bottomrule
  \end{tabular}
  \caption{Diamond-chain adversarial benchmark. Yen-NSP must enumerate
    all $2^{k}$ shortest paths and was halted after exceeding the
    $30$~s budget at $k=15$ (recorded budget overflow:
    $56$~s). CWZ stays in the tens of microseconds across the full
    sweep. The two algorithms are timed in separate passes and each CWZ
    figure is the median of $11$ repetitions; see the measurement note
    below.}
  \label{tab:adv}
\end{table}

\begin{figure}
  \centering
  \includegraphics[width=0.8\textwidth]{../results/adversarial.pdf}
  \caption{Adversarial diamond-chain: Yen-NSP's runtime grows by roughly a
    factor of five per increment of $k$ while CWZ stays microsecond-fast.
    The number of shortest paths Yen must enumerate doubles with $k$, but its
    spur loop rescans the whole result list for every spur index, so the
    measured growth is closer to $4^{k}$ than $2^{k}$.}
  \label{fig:adversarial}
\end{figure}

Between $k=10$ and $k=15$, Yen's runtime grows from $0.023$~s to
$56$~s -- consistent with the $2^{k}$ scaling we expected: each
increment of $k$ doubles the number of shortest paths Yen has to
enumerate before finding a strictly longer one. CWZ, in contrast, finds
the NSP in microseconds across the full sweep because the graph contains
exactly one back-edge (the $s\to t$ shortcut). That edge goes forward in
distance, so \textsc{NextSP-Straight} removes it and records it as a
single candidate (Section~\ref{sec:pipeline}); the answer is found in
constant work after the initial Dijkstra, with the $6$-tuple enumeration
never firing.

At $k=15$ the ratio is
$56.18~\text{s} / 22~\text{\textmu s} \approx 2.6$~million$\times$ in
favour of CWZ. Past $k=15$ Yen is too slow for our budget; CWZ continues
to run unimpeded up to $k=24$ ($|V|=73$) in tens of microseconds.

\paragraph{A measurement note.} An earlier version of this benchmark
interleaved the two algorithms, timing CWZ at each $k$ immediately after
the Yen run at the same $k$. That inflated the CWZ timings at
$k=13,14,15$ by $15$--$50\times$ (to $0.4$--$1.4$~ms), because each CWZ
run began in the allocator and page-table state left behind by a Yen run
that had just churned $2^{k}$ paths. The symptom was a CWZ column that
rose with $k$ up to $k=15$ and then collapsed by a factor of $50$ at
$k=16$, where Yen no longer ran. The harness now measures the two
algorithms in separate passes, with a warm-up and $11$ repetitions per
CWZ point; the resulting CWZ column is flat-to-mildly-increasing across
the whole sweep, as the algorithm's structure predicts. The correction
moves the comparison further in CWZ's favour, but the original figures
were an artefact of the harness rather than of either algorithm.

\paragraph{Significance.} This is the empirical fact CWZ's worst-case
polynomial bound is buying. On any input where the simple heuristic
blows up, the polynomial-time algorithm continues to terminate quickly.
The \emph{benign-random} comparison in the previous sections must be
read in light of this: Yen is faster in practice on random inputs
because the adversarial structure does not arise by accident, but it
cannot be relied on as a worst-case solver. The diamond-chain panel of
Figure~\ref{fig:runtime} shows the same reversal in the routine
benchmark, side-by-side with the four benign families.

\section{Threats to validity}

\paragraph{Fidelity of the implementation.}
\begin{itemize}
  \item Our 2-VDP-in-DAG subroutine is a max-flow-based implementation
    rather than Tholey's linear-time
    algorithm~\cite{Tholey2012}. Both have the same asymptotic
    $O(|V|+|E|)$ bound, but the constant factor differs; we have not
    measured the gap.
  \item Our reduction to $(s,t)$-straight (\texttt{reduce\_to\_straight}
    in \texttt{src/cwz/reductions.cpp}) caps each multi-edge family at
    $K=5$ distinct weights per (source, destination) pair after each
    fold. This is the one place where our implementation is an
    empirically-validated approximation rather than a proven
    transcription: the paper's \textsc{NextSP} instead records
    candidates for expensive routes through a folded vertex and keeps
    only the minimum shortcut. Ablating $K$ against the brute-force oracle
    over $10{,}000$ random instances at $|V|\le 12$ ($7{,}138$ of which
    have an NSP) puts $K=1$ at $2{,}092$ misses ($29.3\%$) and
    $K=2$, $K=3$ and $K=5$ at none, so the margin we actually rely on is
    one extra weight, not five. But no
    constant $K$ is provable: a route can be unusable because the folded
    vertices it passes through are already used elsewhere on the path,
    and nothing bounds how many cheap alternatives such a collision can
    rule out. A pathological input whose only NSP needs the
    $6^{\text{th}}$ or later cheapest shortcut weight cannot be excluded
    a priori.
  \item The layered solver (\texttt{cwz\_nsp\_layered}) is a direct
    transcription of the paper's \textsc{NextSP-Layered}: the single
    $6$-tuple enumeration of Lemma~5.3, on a strictly $(s,t)$-layered
    graph, with the removed-back-edge candidates supplied by the
    reduction. This enumeration \emph{alone} reproduces the brute-force
    oracle on all $34{,}000$ trials.
  \item Both oracles are our own code. The brute-force search shares no
    modules with the CWZ pipeline -- it carries its own Dijkstra
    specifically so that a bug in \texttt{src/shortest\_path/} cannot
    hide inside the oracle that is supposed to catch it -- but it was
    written by one author from one reading of the problem statement. A
    shared misunderstanding of the NSP definition would be invisible to
    both. Yen, implemented independently from a published algorithm, is
    the partial mitigation, and it agrees on every benign instance.
\end{itemize}

\paragraph{Validity of the measurements.}
\begin{itemize}
  \item Runtimes are not reproducible to better than roughly a factor of
    three. The diamond-chain family ignores the random seed, so its
    $|V|=50$ instance is deterministic -- the same graph on every run --
    yet Yen measured $1.54$~s on one run of the sweep and $4.82$~s on
    another, the latter confirmed by five standalone repetitions on an
    idle machine. The main sweep takes a \emph{single} wall-clock
    measurement per instance, with no repetition, on a shared laptop
    under WSL2; only the adversarial sweep of
    Section~\ref{sec:adversarial} repeats each point. No runtime claim in
    this chapter finer than ``orders of magnitude'' should be relied on.
  \item Sub-microsecond runtimes fall below our timer resolution and are
    clamped to a $1$~\textmu s floor for display. The flat segments at
    that level in Figure~\ref{fig:runtime} mean ``too fast to measure'',
    not ``constant time'', and no comparison between algorithms is
    meaningful there.
  \item Twenty instances per cell is a moderate sample: a single instance
    moves a rate by $5$ percentage points. The small dip in
    Figure~\ref{fig:existence} at $|V|=30$ is one such instance and is
    noise -- at $200$ trials the same configuration gives $100\%$. The
    IQR bands on Figure~\ref{fig:runtime} expose the corresponding
    spread in runtime; on Erd\H{o}s--R\'enyi at large $|V|$ the band
    covers roughly a factor of two around the median.
\end{itemize}

\paragraph{Generality of the inputs.}
\begin{itemize}
  \item The sampled distribution is \emph{conditioned}. We redraw any
    instance in which $t$ is unreachable from $s$
    (Section~\ref{sec:methodology}), so what we measure is not
    Erd\H{o}s--R\'enyi $G_{n,p}$ but $G_{n,p}$ conditioned on $s\to t$
    connectivity. At small $|V|$ the two differ substantially -- at
    $|V|=6$ roughly two thirds of raw draws were being rejected. The
    conditioning is the right choice for measuring an NSP solver, since
    a graph with no $s\to t$ path poses no NSP question at all, but our
    results describe the conditioned distribution and do not transfer
    unchanged to the unconditioned one.
  \item The horizontal axis is the \emph{target} $|V|$, and the actual
    size differs by family: at target $50$ the grid instances have $64$
    vertices, the layered instances $51$, and the diamond chains $49$.
    Comparisons \emph{within} a family across sizes are sound;
    comparisons \emph{between} families at a fixed abscissa are not
    strictly like-for-like.
  \item The generator families we sweep -- excluding diamond-chain --
    are benign: none deliberately constructs many alternate shortest
    paths to stress Yen. The dedicated diamond-chain sweep
    (Section~\ref{sec:adversarial}) addresses this gap directly, but it
    is a single hand-built construction rather than a family of
    adversarial structures.
  \item Yen's enumeration cap $k_{\max}=5000$ is a configuration choice,
    and it is what fixes the point at which Yen stops merely being slow
    and starts returning the wrong answer (Section~\ref{sec:agreement}).
    That a capped Yen must eventually fail on graphs with exponentially
    many shortest paths is unavoidable; that it fails specifically at
    $|V|\ge 40$ on this family is an artefact of the cap we chose.
  \item All experiments stop at $|V|=50$ and use integer weights drawn
    uniformly from $[1,10]$. Narrow weight ranges make ties -- and hence
    equal-cost shortest paths -- more likely than a continuous or wider
    distribution would, which favours neither algorithm uniformly: ties
    are what defeat both Yen and the brute-force bound, while CWZ is
    indifferent to them.
\end{itemize}

\chapter{Conclusion}
\label{ch:conclusion}

This thesis presented the first implementation and empirical study of
the Chen--Wein--Zhang next-to-shortest-path algorithm. The algorithm
resolves a thirty-year-old open question for directed graphs with
strictly positive edge weights and stands out as the only NSP algorithm
with a fixed polynomial bound on general inputs.

\section{Summary of findings}

\begin{enumerate}
  \item The algorithm is implementable from the paper's pseudocode
    without needing access to optimised 2-VDP subroutines: a
    straightforward max-flow-based 2-VDP yields a working implementation
    at the cost of a higher constant factor (the asymptotic
    $O(|V|+|E|)$ bound is the same as Tholey's linear-time algorithm).
  \item On the four benign random graph families we measured
    (Erd\H{o}s--R\'enyi, random DAG, layered, grid) at $|V|\le 50$,
    Yen's $k$-shortest-simple-paths algorithm is several orders of
    magnitude faster than CWZ. This was expected: our random generators
    almost never produce many distinct simple shortest paths, so Yen
    rarely takes more than a handful of iterations.
  \item On a deliberate adversarial family of \emph{diamond-chain}
    graphs designed to have $2^{k}$ distinct shortest $s\to t$ paths
    (Section~\ref{sec:adversarial}), the picture reverses dramatically.
    At $k=15$ ($|V|=46$, $|E|=61$) Yen takes $59$~seconds while CWZ
    finishes in $1.4$~milliseconds --- a ${\sim}41{,}000\times$ ratio
    that grows exponentially with $k$. CWZ scales smoothly to $k=24$
    ($|V|=73$) in tens of microseconds; Yen would require many hours on
    the same instance by extrapolation. Worse, with a practical
    enumeration cap Yen does not merely slow down: at $|V|\ge 40$ it
    exceeds its cap and returns the \emph{wrong} answer (``no NSP'')
    while CWZ stays correct. This is the empirical demonstration of why
    a polynomial-time NSP algorithm matters.
  \item Our implementation is a faithful transcription of the paper's
    three-stage algorithm: \textsc{NextSP} (fold off-shortest-path
    vertices), \textsc{NextSP-Straight} (remove forward/sideways
    back-edges, recording their candidate paths, to obtain a strictly
    $(s,t)$-layered graph), and \textsc{NextSP-Layered} (the Lemma~5.3
    $6$-tuple enumeration). It matches the brute-force oracle with
    \emph{zero} mismatches across more than $35{,}000$ random instances
    ($|V|\le 14$) and matches Yen on every benign benchmark instance up
    to $|V|=50$. An ablation confirms the $6$-tuple enumeration alone
    suffices on the strictly-layered graph --- the property the paper's
    correctness proof relies on. The one remaining approximation is the
    $K=5$ multi-edge cap in \textsc{NextSP} (below).
\end{enumerate}

\section{Limitations and future work}

\paragraph{Engineering improvements.} Replacing the max-flow 2-VDP by
Tholey's linear-time algorithm~\cite{Tholey2012} would lower the
constant factor on the inner subroutine. Eliminating the
\textsc{NextSP-Straight} subdivision via clever distance indexing,
suggested by the paper's authors as a clarity-oriented choice, could
provide further speedup.

\paragraph{Larger graphs.} Our experiments stop at $|V|=50$. The
empirical curve on Erd\H{o}s--R\'enyi suggests the actual runtime grows
polynomially with an exponent of roughly $3$, well below the
worst-case $|V|^{4}|E|^{3}$ bound; characterising this gap on still
larger or differently-structured inputs would be valuable.

\paragraph{Multi-edge cap tightening.} The $K=5$ multi-edge cap in
\texttt{reduce\_to\_straight} is the one approximation in our pipeline.
The faithful fix is to implement \textsc{NextSP}'s own
candidate-recording step (mirroring what \textsc{NextSP-Straight}
already does for back-edges in our code): record the candidate path
through each folded vertex, then keep only the single minimum-weight
shortcut per pair. This would eliminate both the cap and the Cartesian
multi-edge blow-up that motivated it, making the entire pipeline a
proven transcription of the paper.

\paragraph{Adversarial benchmarks.} The diamond-chain family is the
simplest construction we know of that defeats Yen. Richer adversarial
graphs --- recursive nestings, irregular branching factors, weight
ratios that vary with depth --- would map out the boundary between
``Yen survives'' and ``CWZ is the only viable option'' more finely.

\paragraph{Algorithmic refinements.} The paper's authors note their
bound is ``not optimised''. Candidate optimisations include exploiting
layer structure to prune the $6$-tuple enumeration and amortising 2-VDP
queries across nearby tuples, as well as replacing our max-flow 2-VDP
with Tholey's linear-time routine.


\bibliographystyle{plain}
\bibliography{refs}

\end{document}
